% Bagian: sets-functions-relations
% Bab: relations
% Subbagian: orders

\documentclass[../../../include/open-logic-section]{subfiles}

\begin{document}

\olfileid[id]{sfr}{rel}{ord}
\olsection{Urutan}

\begin{explain}
Banyak perbandingan kita melibatkan penyifatan objek tertentu sebagai
``lebih kecil daripada'', ``sama dengan'', atau ``lebih besar daripada'' objek
lain dalam segi tertentu. Perbandingan ini melibatkan relasi \emph{urutan}.
Namun, terdapat berbagai jenis relasi urutan. Sebagai contoh, sebagian relasi
mengharuskan setiap dua objek dapat dibandingkan, sedangkan relasi lain tidak.
Sebagian mencakup identitas (seperti~$\le$), dan sebagian mengecualikannya
(seperti~$<$). Taksonomi berikut akan membantu kita.
\end{explain}

\begin{defn}[Praurutan]
Suatu relasi yang sekaligus refleksif dan transitif disebut
\emph{praurutan.}
\end{defn}

\begin{defn}[Urutan parsial]
Suatu praurutan yang juga antisimetris disebut \emph{urutan parsial}.
\end{defn}

\begin{defn}[Urutan linear]\ollabel{def:linearorder}
Suatu urutan parsial yang juga terhubung disebut \emph{urutan total} atau
\emph{urutan linear.}
\end{defn}

\begin{ex}
Setiap urutan linear juga merupakan urutan parsial, dan setiap urutan parsial
juga merupakan praurutan, tetapi implikasi sebaliknya tidak berlaku. Relasi
universal pada~$A$ merupakan praurutan karena relasi itu refleksif dan
transitif. Namun, jika $A$ mempunyai lebih dari satu !!{element}, relasi
universal tidak antisimetris sehingga bukan urutan parsial.
\end{ex}

\begin{ex}
Tinjaulah relasi \emph{tidak lebih panjang daripada} $\preccurlyeq$
pada~$\Bin^*$: $x \preccurlyeq y$ jika dan hanya jika $\len{x} \le \len{y}$.
Relasi ini merupakan praurutan (refleksif dan transitif), bahkan terhubung,
tetapi bukan urutan parsial karena tidak antisimetris. Sebagai contoh,
$01 \preccurlyeq 10$ dan $10 \preccurlyeq 01$, tetapi $01 \neq 10$.
\end{ex}

\begin{ex}
Salah satu urutan parsial yang penting adalah relasi $\subseteq$ pada suatu
himpunan yang beranggotakan himpunan-himpunan. Relasi ini pada umumnya bukan
urutan linear. Jika $a \neq b$ dan kita meninjau
$\Pow{\{a, b\}} = \{\emptyset, \{a\}, \{b\}, \{a,b\}\}$, tampak bahwa
$\{a\} \nsubseteq \{b\}$, $\{a\} \neq \{b\}$, dan
$\{b\} \nsubseteq \{a\}$.
\end{ex}

\begin{ex}
Relasi \emph{keterbagian tanpa sisa} memberi kita suatu urutan parsial yang
bukan urutan linear. Untuk bilangan bulat $n$ dan~$m$, kita menulis $n \mid m$
untuk menyatakan bahwa $n$ membagi habis $m$, yaitu jika dan hanya jika ada
suatu bilangan bulat~$k$ sehingga $m = kn$. Pada~$\Nat$, relasi ini merupakan
urutan parsial, tetapi bukan urutan linear; misalnya, $2 \nmid 3$ dan juga
$3 \nmid 2$. Jika dipandang sebagai relasi pada $\Int$, keterbagian hanya
merupakan praurutan karena tidak antisimetris: $1 \mid -1$ dan $-1 \mid 1$,
tetapi $1 \neq -1$.
\end{ex}

\begin{ex}
\emph{Relasi perluasan} pada suatu himpunan barisan~$A^*$ adalah sebagai
berikut: $s \sqsubseteq s'$ jika dan hanya jika $s = \emptyseq$ (barisan
kosong), $s = s'$, atau $s = \tuple{s_1, \dots, s_n}$ dan
$s' = \tuple{s_1, \dots, s_n, s_{n+1}, \dots, s_m}$. Jika
$s \sqsubseteq s'$, kita juga mengatakan bahwa $s$ merupakan \emph{segmen
awal} dari~$s'$. Relasi perluasan pada $A^*$ merupakan urutan parsial, tetapi
bukan urutan linear. Sebagai contoh, jika $a \neq b$, maka
$ab \not\sqsubseteq ba$ dan $ba \not\sqsubseteq ab$.
\end{ex}

\begin{defn}[Urutan ketat]
Suatu \emph{urutan ketat} adalah relasi yang irefleksif, asimetris, dan
transitif.
\end{defn}

\begin{defn}[Urutan linear ketat]\ollabel{def:strictlinearorder}
Suatu urutan ketat yang juga terhubung disebut \emph{urutan total ketat} atau
\emph{urutan linear ketat.}
\end{defn}

\begin{ex}
$\le$ merupakan urutan linear yang bersesuaian dengan urutan linear
ketat~$<$. $\subseteq$ merupakan urutan parsial yang bersesuaian dengan
urutan ketat~$\subsetneq$.
\end{ex}

Setiap urutan ketat $R$ pada~$A$ dapat diubah menjadi urutan parsial dengan
menambahkan diagonal $\Id{A}$, yakni dengan menambahkan semua pasangan
$\tuple{x, x}$. (Ini disebut \emph{penutupan refleksif} dari~$R$.) Sebaliknya,
dengan memulai dari suatu urutan parsial, kita dapat memperoleh urutan ketat
dengan menghapus~$\Id{A}$. Dua hasil berikut memperjelas pernyataan ini.

\begin{prop}\ollabel{prop:stricttopartial}
Jika $R$ merupakan urutan ketat pada~$A$, maka $S = R \cup \Id{A}$ merupakan
urutan parsial. Selain itu, jika $R$ merupakan urutan linear ketat, maka $S$
merupakan urutan linear.
\end{prop}

\begin{proof}
Andaikan $R$ merupakan urutan ketat, yakni $R \subseteq A^2$ serta $R$
irefleksif, asimetris, dan transitif. Misalkan $S = R \cup \Id{A}$. Kita
harus menunjukkan bahwa $S$ refleksif, antisimetris, dan transitif.

$S$ jelas refleksif karena $\tuple{x, x} \in \Id{A} \subseteq S$ untuk
semua $x \in A$.

Untuk menunjukkan bahwa $S$ antisimetris, andaikan demi reductio bahwa
$Sxy$ dan $Syx$, tetapi $x \neq y$. Karena
$\tuple{x,y} \in R \cup \Id{A}$, sedangkan
$\tuple{x, y} \notin \Id{A}$, harus berlaku $\tuple{x, y} \in R$, yakni
$Rxy$. Dengan cara yang sama, berlaku~$Ryx$. Namun, hal ini bertentangan
dengan asumsi bahwa $R$ asimetris.

Untuk transitivitas, andaikan $Sxy$ dan $Syz$. Jika
$\tuple{x, y} \in R$ dan $\tuple{y,z} \in R$, maka
$\tuple{x, z} \in R$ karena $R$ transitif. Jika tidak, maka
$\tuple{x, y} \in \Id{A}$, yakni $x = y$, atau
$\tuple{y, z} \in \Id{A}$, yakni $y = z$. Dalam kasus pertama, menurut asumsi
berlaku $Syz$ dan $x = y$, sehingga $Sxz$. Kasus kedua serupa. Dalam kedua
kasus, $Sxz$; jadi, $S$ juga transitif.

Untuk klausa ``selain itu'', andaikan bahwa $R$ juga terhubung. Jadi, untuk
semua $x \neq y$, berlaku $Rxy$ atau~$Ryx$, yakni
$\tuple{x, y} \in R$ atau $\tuple{y, x} \in R$. Karena $R \subseteq S$,
pernyataan ini tetap berlaku untuk $S$; jadi, $S$ juga terhubung.
\end{proof}

\begin{prop}\ollabel{prop:partialtostrict}
Jika $R$ merupakan urutan parsial pada~$A$, maka
$R^- = R \setminus \Id{A}$ merupakan urutan ketat. Selain itu, jika $R$
merupakan urutan linear, maka $R^-$ merupakan urutan linear ketat.
\end{prop}

\begin{proof}
Hal ini diserahkan sebagai latihan.
\end{proof}

\begin{prob}
Berikan bukti untuk \olref[sfr][rel][ord]{prop:partialtostrict}.
\end{prob}

Hasil sederhana berikut menetapkan bahwa urutan linear ketat memenuhi sifat
yang mirip dengan ekstensionalitas.

\begin{prop}\ollabel{prop:extensionality-strictlinearorders}
Jika $<$ merupakan urutan linear ketat pada $A$, maka:
\[
  (\forall a, b \in A)((\forall x \in A)(x < a \liff x < b) \lif a = b).
\]
\end{prop}

\begin{proof}
Andaikan $(\forall x \in A)(x < a \liff x < b)$. Jika $a < b$, maka $a < a$,
bertentangan dengan fakta bahwa $<$ irefleksif; jadi, $a \nless b$. Dengan
cara yang persis sama, $b \nless a$. Jadi, $a = b$, karena $<$ terhubung.
\end{proof}

\end{document}
